% chapters/ch02f_state_labels.tex
% Section 2.9: State Labels (HST-N to SOP Interface)

\section{State Labels (HST-N to SOP Interface)}
\label{sec:state-labels-c2}

\nrmobox{HST-N non-operational charter}{
HST-N is the \emph{state-topology map} for NRMO. It is not an
operational module; it produces only state classifications. SOP is
the consumer interface; HST-N never executes.
}

\subsection{HST-N State Labels (SOP Input)}
\label{sec:hstn-state-labels-c2}

\begin{table}[ht]
\centering
\small
\begin{tabularx}{\linewidth}{lXX}
\toprule
\textbf{Label} & \textbf{Meaning} & \textbf{Design intent} \\
\midrule
LOAD
& Intensity of load where judgement errors increase
& When errors increase: accumulate reversible trials. \\
VOL
& Situational volatility (instability)
& When volatility is high: refrain from boundary-approaching
actions. \\
IRREV
& Proximity to irreversible boundary
& Category $D$ requires procedure; may require SHUTDOWN. \\
\bottomrule
\end{tabularx}
\caption[HST-N State Labels (SOP Input)]{HST-N State Labels (Table
2.4). Labels are produced by HST-N and consumed by the SOP layer.}
\label{tab:hstn-state-labels-c2}
\end{table}

\subsection{State Label Value Definitions}
\label{sec:hstn-state-label-values-c2}

\begin{table}[ht]
\centering
\small
\begin{tabularx}{\linewidth}{llX}
\toprule
\textbf{Label} & \textbf{Values} & \textbf{Description} \\
\midrule
LOAD
& $L_0$ / $L_1$ / $L_2$
& $L_0$ = low load, $L_1$ = medium load, $L_2$ = high load.
Higher $\Rightarrow$ more judgement / dialogue errors likely. \\
VOL
& $V_0$ / $V_1$ / $V_2$
& $V_0$ = stable, $V_1$ = volatile, $V_2$ = highly volatile.
Higher $\Rightarrow$ situation more likely to break down. \\
IRREV
& $I_0$ / $I_1$
& $I_0$ = boundary distant, $I_1$ = boundary near.
Near $\Rightarrow$ higher irreversible-transition risk. \\
\bottomrule
\end{tabularx}
\caption[State Label Value Definitions]{State Label Value
Definitions (Table 2.5). Each label takes one value at any given
time.}
\label{tab:hstn-state-label-values-c2}
\end{table}

\subsection{Observation Inputs}
\label{sec:hstn-observation-inputs-c2}

HST-N classifies observations into the discrete state labels. The
classification is non-judgemental.

\paragraph{Observation classes.}

\begin{itemize}
\item Subjective: mood self-report, fatigue self-report.
\item Objective: sleep duration, work hours, recent commitments,
financial state, relational signals.
\item Behavioural: deviation from baseline patterns; blame-shifting
behaviour at boundaries; out-of-band communication.
\end{itemize}

\subsection{Output Specification (SOP Output)}
\label{sec:hstn-output-spec-c2}

The HST-N output is consumed by SOP only:

\begin{itemize}
\item LOAD --- load-judgement signal (when high $\Rightarrow$
restrict to reversible).
\item VOL --- volatility signal (when $V_2$ $\Rightarrow$ refrain
from boundary actions).
\item IRREV --- irreversibility-imminent signal (when $I_1$
$\Rightarrow$ category $D$ requires procedure or SHUTDOWN).
\end{itemize}

\nrmobox{HST-N principle}{
HST-N \textbf{does not propose actions}. Observations $\to$
classifications $\to$ SOP input. SOP retains all decision authority.
}

\subsection{Irreversibility Types}
\label{sec:hstn-irreversible-types-c2}

The irreversibility boundary is not one-dimensional. HST-N
classifies four types of irreversibility:

\begin{table}[ht]
\centering
\small
\begin{tabularx}{\linewidth}{lXXX}
\toprule
\textbf{Type} & \textbf{What decreases} & \textbf{Common triggers} & \textbf{Avoidance principle} \\
\midrule
R-Trust
& Reputation, relational credibility
& Public statements, broken commitments
& Slow tempo; reversible expressions; confirmation. \\

R-Resource
& Capital (time, money, attention)
& Over-commitment, ultimatum-driven binding
& Limit-based commitment; staged release. \\

R-Relationship
& Specific bilateral relationship integrity
& Ultimatums, public airing of grievances
& Slow tempo; private confirmation; no broadcast. \\

R-Agency
& Personal autonomy and self-direction
& Coercion, dependency formation
& Confirm own action; third-party check. \\
\bottomrule
\end{tabularx}
\caption[Irreversibility types]{Four types of irreversibility used
by HST-N. Each type has its own avoidance principle, applied by
SOP based on HST-N classification.}
\label{tab:hstn-irreversibility-types-c2}
\end{table}

\paragraph{NRMO interpretation of irreversibility.}
NRMO treats irreversibility as a \emph{first-class loss}: the
contraction of future degrees of freedom. The avoidance principle
above is the SOP-level operationalisation.

\subsection{HST-N Boundary Principle}
\label{sec:hstn-boundary-principle-c2}

HST-N never crosses the boundary into evaluation or persuasion.

\begin{itemize}
\item HST-N does \emph{not} say ``the situation is good'' or ``the
situation is bad.''
\item HST-N's outputs are classifications, never evaluations.
\item Strong observations are passed to SOP without
interpretation.
\item Output flow is one-way pass-through: HST-N $\to$ SOP. No
direct action.
\end{itemize}

\subsection{HST-N Operational Document File}
\label{sec:hstn-html-c2}

The reference HST-N implementation is locally generated and stored
in the file \texttt{hst-n.html}. The HTML file restates the
specification in form-fillable shape: NRMO OS state premise,
permitted-state transitions, and audit log fields.

\subsection{Operational Declaration (Non-Operational)}
\label{sec:hstn-operational-declaration-c2}

\nrmobox{Non-Operational Charter (restated)}{
HST-N maps state. HST-N does not change state.\\[0.3em]
HST-N's outputs are SOP inputs. HST-N is never standalone.
}

\paragraph{Purpose.}
The HST-N charter is to provide a stable, non-judgemental,
non-operational mapping of subject state. It exists to keep SOP
inputs clean.

\paragraph{Relation to NRMO.}
NRMO treats the HST-N output as factual classification. NRMO does
not negotiate with HST-N classifications: if HST-N reports $L_2$,
NRMO assumes $L_2$.

\subsection{NRMO Layered Diagram}
\label{sec:hstn-nrmo-layered-diagram-c2}

\begin{lstlisting}[basicstyle=\ttfamily\small,frame=single]
[ NRMO (OS / governance / final authority) ]
                    ^
                    | (veto only)
                    |
[ Observation OS / monitoring / classification ]
                    ^
                    |
              HST-N labels
                  (LOAD, VOL, IRREV)
\end{lstlisting}

\paragraph{NRMO authority.}
NRMO retains final veto. HST-N feeds; NRMO decides; Type ZERO
operationalises.

\subsection{Permitted / Prohibited (Restated)}
\label{sec:hstn-permitted-prohibited-c2}

\paragraph{Permitted (HST-N may).}

\begin{itemize}
\item Observe, classify, log.
\item Publish negative assertions.
\item Publish Red Flag codes.
\item Publish Hold notifications.
\end{itemize}

\paragraph{Prohibited (HST-N must not).}

\begin{itemize}
\item Make judgements (good/bad).
\item Cross irreversibility-boundary commentary
(``you should / must / shouldn't'').
\item Issue actions (NRMO and Type ZERO and SOP own this).
\item Modify own classification thresholds without explicit external
re-approval.
\end{itemize}

\subsection{Final Declaration}
\label{sec:hstn-final-declaration-c2}

\begin{lstlisting}[basicstyle=\ttfamily\small,frame=single]
STATE = DAILY_NRMO
\end{lstlisting}

\paragraph{Meaning.}
The default operational state is daily NRMO. HST-N classifications
deviate the state only when their thresholds are crossed; the
default is a no-action posture.
